\documentclass[10pt,a4paper]{amsart}
\usepackage[margin=19mm]{geometry}
\usepackage{amsmath,amssymb,mathtools}
\usepackage[scr=rsfs,cal=euler]{mathalfa}
\usepackage{xfrac}
\usepackage[unicode,hidelinks,bookmarks=false]{hyperref}
\newcommand{\ie}{i.e.\ }
\newcommand{\cf}{cf.\ }
\newcommand{\resp}{respectively\ }
\newcommand{\Subsection}{\subsection}
\providecommand{\nobreakdash}{\nobreak\hbox{-}}
\setlength{\emergencystretch}{2em}
\multlinegap0pt
\usepackage{bm}
\usepackage{bbm}
\usepackage{dsfont}
\usepackage{xspace}
\usepackage{cancel}
\usepackage{mathtools}
\usepackage{amsthm}
\makeatletter
\@ifundefined{theorem}{\newtheorem{theorem}{Theorem}}{}
\@ifundefined{lemma}{\newtheorem{lemma}[theorem]{Lemma}}{}
\@ifundefined{proposition}{\newtheorem{proposition}[theorem]{Proposition}}{}
\makeatother
\title[A Generalisation on Erd\H{o}s Distinct Subset Sums Problem]{A Generalisation on Erd\H{o}s Distinct Subset Sums Problem}
\author{Zijie Gu}
\date{Source: arXiv:2510.06032v1 (CC BY 4.0)}
\begin{document}
\maketitle

\noindent\emph{Source and licence.} A Generalisation on Erd\H{o}s Distinct Subset Sums Problem. Version arXiv:2510.06032v1 (CC BY 4.0), distributed under the Creative Commons Attribution 4.0 International licence, as recorded in the arXiv licence metadata for this exact version and shown on the abstract page https://arxiv.org/abs/2510.06032v1. Full rights notice: licenses/arxiv-2510.06032-CC-BY-4.0.txt.\par\smallskip

\noindent\textit{Editorial scope.} Counted unit: the Theorem whose source label is \texttt{thm:first} in the article (source0.tex lines 165--170, source label \texttt{thm:first}), delivered together with the article's own complete original proof of it (source0.tex lines 172--202, one \texttt{proof} environment of 31 lines). No mathematics is written, repaired or re-derived: statement and proof are the article's own text line for line. Required context retained uncounted: thm:pnorm-volume (lines 120--130); lem:pnorm (lines 139--144). Every definition, display and label of those lines is retained unchanged. Omitted from this package and not redistributed: 1--119, 131--138, 145--164, 171--171, 203--286. Nothing inside a retained excerpt is omitted or altered: every retained body is delivered exactly as the article's lines are written, so no omission receipt of any kind accompanies this package. Citations: the retained excerpts carry no citation command at all, so no bibliography entry of the article is transcribed and no reference is redistributed. Reference adaptation: none was needed and none was made; every reference inside the delivered unit resolves inside it, so no unresolved reference or citation is produced. Printed slips kept literal: no printed word, symbol, hypothesis or number of the counted statement or of its proof was altered. Layout and class: the article's own class is \texttt{amsart} with packages including amsmath, amssymb, amsthm, graphicx, hyperref, geometry, cite; the wrapper is the lane's pinned \texttt{amsart} editorial wrapper at 10pt on A4, which restates the theorem-like environments the retained text needs and adds the article's own macro definitions that the retained text needs (none), with extra wrapper packages (bm, bbm, dsfont, xspace, cancel, mathtools). The delivered numbering is the unit's own. Assets: no retained span contains a figure environment, a graphics-inclusion command, a PDF-page inclusion, a table or a TikZ picture; no image file is copied into this package, the package declares no asset dependency and no image file is redistributed with it. Licence: version arXiv:2510.06032v1 of this article is distributed under the Creative Commons Attribution 4.0 International licence, as recorded in the arXiv licence metadata for this exact version and shown on the abstract page https://arxiv.org/abs/2510.06032v1. A full rights notice is delivered at licenses/arxiv-2510.06032-CC-BY-4.0.txt. Characters: every retained excerpt is pure ASCII, so the delivered engine, XeLaTeX with its default Latin Modern text font, reports no missing character.
   \begin{proposition}\label{thm:pnorm-volume}
   The volume $V_{k,p}(R)$ of a $k$-dimensional ball with radius $R$ in the $p$-norm, that is,
   \[
   B_{k,p}(R) = \left\{ X \in \mathbb{R}^k : \|X\|_p \leq R \right\},
   \]
   is given by
   \[
   V_{k,p}(R) = \frac{\left[2\Gamma\left(1+\frac{1}{p}\right)\right]^k}{\Gamma\left(1+\frac{k}{p}\right)} \cdot R^k,
   \]
   where $\Gamma$ denotes the Gamma function.
   \end{proposition}

   \begin{lemma}\label{lem:pnorm}
   In a $p$-normed space, let $A$ denote the set of $2^n$ integer lattice points (i.e., points in $\mathbb{Z}^k$) that are closest to the origin. Let $R$ be the radius of the smallest ball (centred at the origin) that contains all points in $A$. Then the following identity holds:
   \[
   \sum_{s\in A}{\|s\|_p^p}=2^n\cdot \frac{k}{k+p}\cdot R^p
   \]
   \end{lemma}

\begin{theorem}\label{thm:first}
We have:
\[
M \;\ge\; (1+o(1))\;\cdot\;\frac{1}{k+1}\;\cdot\;\sqrt{\frac{\pi}{2}}\;\cdot\;\frac{2^{\,\tfrac{n}{k}}}{\sqrt{n}}\;\cdot\;(k!)^{1/k}.
\]
\end{theorem}

\begin{proof}
Since the $k$ components of the vector $X$ are independent of each other, we have:
\[
\mathbb{E}[\|X\|_1] \leq k \cdot \max \mathbb{E}\left[\left|\sum_{i=1}^n \epsilon_i x_i\right|\right],
\]
where each $x_i$ is between $0$ and $M$.

The function $\mathbb{E}\left[\left|\sum_{i=1}^n \epsilon_i x_i\right|\right]$ is convex with respect to each variable $x_i$. Therefore, its maximum is achieved when each $x_i$ takes either the value $0$ or $M$. Let $n_0$ denote the number of $x_i$ equal to $M$, and the remaining $x_i$ are $0$. In this case, we have:
\begin{align*}
    \mathbb{E}\left[\left|\sum_{i=1}^n \epsilon_i x_i\right|\right] &= \frac{1}{2^n} \sum_{\epsilon_i} \left|\sum_{i=1}^n \epsilon_i x_i\right| \\
    &\leq \frac{1}{2^n} \cdot 2^{n-n_0} \cdot M \cdot \sum_{\epsilon_i} \left|\sum_{i=1}^{n_0} \epsilon_i\right| \\
    &= \frac{M}{2^{n_0}} \sum_{i=0}^{n_0} \binom{n_0}{i} \left| \frac{n_0}{2} - i \right| \\
    &= \frac{M}{2^{n_0}} \cdot n_0 \cdot \binom{n_0-1}{\left\lfloor \frac{n_0-1}{2} \right\rfloor} \\
    &\leq \frac{M}{2^n} \cdot n \cdot \binom{n-1}{\left\lfloor \frac{n-1}{2} \right\rfloor}
\end{align*}

Therefore,
\begin{equation}
\mathbb{E}[|X|]\le \frac{kMn}{2^n}\cdot \binom{n-1}{\lfloor\frac{n-1}{2}\rfloor}
\end{equation}

On the other hand, Lemma~\ref{lem:pnorm} (with $p=1$) gives a lower bound:
\begin{equation}
\mathbb{E}[|X|] \geq (1+o(1)) \frac{k}{k+1} R,
\end{equation}
where (using Proposition~\ref{thm:pnorm-volume})
\[
R=2^{n/k}\cdot \frac{\Gamma \left(\frac{k+1}{k}\right)^{1/k}}{2\Gamma (2)}=2^{n/k-1}\cdot \Gamma \left(\frac{k+1}{1}\right)^{1/k}=2^{n/k-1}\cdot \left(k!\right)^{1/k}
\]
Comparing these two bounds and solving for $M$ gives the result in the theorem~\ref{thm:first}.
\end{proof}

\end{document}
